% !TEX root = final_report.tex
\documentclass[final_report.tex]{subfiles}
\setcounter{chapter}{1}

\begin{document}
\ourchapter{Background and related work}\label{ch:background}

\section{Self-supervised visual representation learning}\label{sec:ssl}
\stub{The problem setting and the main families --- contrastive, masked
reconstruction, and pretext-task methods --- with enough detail to place the
position-prediction lineage inside it.}

\section{The position-prediction lineage}\label{sec:position-prediction}
\stub{Context prediction and jigsaw puzzles as the ancestors; MP3, DropPos and MAE
as the modern absolute-position and masked-reconstruction methods. Evaluates each
approach rather than merely listing it: what each supervises, and what that costs.}

\section{PART in detail}\label{sec:part}
\stub{The source paper's method. Off-grid continuous patch sampling, pairwise
relative targets, the cross-attention relative encoder, and the absence of position
embeddings. Includes the scale and aspect extension of Section~5.1(ii), which is the
precedent for adding a channel group.}

\section{Rotation in self-supervision}\label{sec:rotation-in-ssl}
\stub{RotNet's discrete image-level rotation prediction, and why \emph{relative}
orientation between patches is a different object. A short contrast with
architecturally rotation-equivariant approaches: they impose invariance, whereas
RoPART treats orientation as recoverable signal to be supervised.}

\section{Orientation-sensitive tasks and their metrics}\label{sec:orientation-tasks}
\stub{Horizon line estimation and head pose estimation --- what makes a task
orientation-sensitive, and how each is scored. Motivates the choice of HLW and the
decision to headline angular error rather than the published AUC.}

\section{Positioning RoPART}\label{sec:positioning}
\stub{The gap: PART names ``modelling rotations'' as open future work. States what
RoPART adds and how it relates to everything above.}

\ifSubfilesClassLoaded{\pagebreak\printbibliography}{}
\end{document}
